\section{Síntesis y ampliación}

\subsection{Un hilo conductor que atraviesa tres líneas de trabajo}

El \textbf{preentrenamiento en dos fases} demuestra que el \textbf{orden} de los datos tiene valor; el \textbf{front-loading del razonamiento} muestra que el \textbf{momento} de los datos relacionados con la capacidad es relevante; y RLP indica que el \textbf{comportamiento y las recompensas} antes de la predicción son importantes. Juntos, estos tres enfoques impulsan el preentrenamiento desde un ajuste estático de máxima verosimilitud hacia un aprendizaje consciente del currículo, consciente de las etapas y consciente de la acción.

\begin{importantbox}{Marco de conocimiento final}
\begin{enumerate}
  \item La mezcla de datos debe registrar simultáneamente los pesos y el cronograma temporal.
  \item Calidad y época/repetición son variables distintas; incluso los datos de alta calidad tienen un límite de repetición.
  \item La asignación temprana de datos de razonamiento puede alterar el fundamento base disponible para la etapa posterior.
  \item RLP mide la ganancia informativa predictiva del pensamiento mediante contrafactuales sin pensar (no-think).
  \item Las recompensas densas, la línea base EMA y la actualización relativa al grupo determinan conjuntamente la estabilidad.
  \item Todos los beneficios deben interpretarse bajo las métricas de tokens, FLOPs, checkpoint, post-entrenamiento y evaluación en benchmarks.
\end{enumerate}
\end{importantbox}

\subsection{De los resultados del currículo a las preguntas de investigación}

En primer lugar, ¿es posible aprender un currículo continuo automáticamente en lugar de hacerlo manualmente en dos fases? En segundo lugar, ¿cómo se retroalimenta el sesgo del clasificador de calidad hacia el comportamiento del modelo? En tercer lugar, ¿es igualmente efectivo el front-loading para el razonamiento no STEM, multilingüe y multimodal? En cuarto lugar, ¿pueden incorporarse restricciones de facticidad, fidelidad causal o seguridad a la recompensa del pensamiento sin comprometer la escalabilidad? Y en quinto lugar, ¿puede reducirse el cómputo del rollout de RLP mediante pensamiento latente, generación especulativa o un batching más robusto?

\begin{knowledgebox}{Experimento factorial sugerido}
Es posible utilizar un diseño $2\times2\times2$ para activar y desactivar dos fases, front-loading y RLP por separado, midiendo los efectos principales y las interacciones bajo una arquitectura fija, un total de tokens constante, FLOPs aproximados y una receta de post-entrenamiento definida. Solo así se podrá responder si las tres estrategias son complementarias, redundantes o mutuamente interferentes, en lugar de ensamblar mecánicamente los resultados de tres artículos diferentes.
\end{knowledgebox}

\subsection{Lecturas adicionales}

Se recomienda leer siguiendo el orden de los problemas: primero, el preentrenamiento en dos fases de Feng et al., para comprender la mezcla de datos y el currículo; luego, el front-loading del razonamiento de Akter et al., verificando las condiciones de reason-base y las cinco lecciones; finalmente, RLP v2 de Hatamizadeh et al., centrándose en la identidad de mejora esperada, el algoritmo 1, el emparejamiento de cómputo y las ablaciones. Posteriormente, contrastar Quiet-STaR, RPT y RLPT comparando sus fuentes de recompensa y su granularidad.

\begin{warningbox}{Mantener la instantánea de la clase al leer los artículos}
Esta clase se impartió el 2026-04-30: dos fases utilizó arXiv v1, front-loading utilizó v1 y RLP utilizó v2 del 2026-03-01. Revisiones posteriores, reproducciones o resultados de producto pueden servir como nueva evidencia, pero no se deben retrotraer para convertirlos en hechos que ya estaban establecidos el día de la clase.
\end{warningbox}

Finalmente, el «inteligencia de próxima generación» propuesta por el curso no es un modelo nuevo aislado, sino una visión del entrenamiento: bajo datos y potencia computacional limitados, los modelos deben elegir con más inteligencia el orden de aprendizaje, establecer más tempranamente un fundamento de razonamiento y permitir que los comportamientos de exploración reciban retroalimentación escalable y comparable. El futuro del preentrenamiento podría no residir únicamente en un $N$ mayor, sino también en un proceso de aprendizaje mejorado.

\end{document}
